\section{楚：江汉以南，另一种天下}

楚在江汉之间形成政治中心时，北方的文献往往已预先给它安放了一个位置：“夷”“荆”或不守周礼的对手。这样的称谓不能替代楚自身的国家史。江汉的河流把平原、上游山道、淮域的城邑和更南面的空间接在一起，也把治理的难题带到楚廷面前：军队要沿水陆通道移动，粮秣要经过地方城邑，国君的命令还须经令尹、公族与受楚保护或压迫的小国才能落实。楚的力量不在于远离中原，而在于能够把这套南方腹地转化为北进的兵员、道路和关系；它的限制，也正藏在转换并不自动完成的地方。

\eventanchor{SA-0704-CHU-KING}{春秋·楚君称王}

传统年表把约前704年熊通自立为楚武王列为一条醒目的界线。此时《春秋》尚未开始编年，因而不能把这个年份误写成鲁国经文已经逐年见证的事实；后来的《史记·楚世家》将它纳入楚国兴起的长线。不过，“王”号所提出的问题并不因此虚假。周天子仍掌握一套可被各国理解的名分，楚却以最高称号拒绝只在北方诸侯所分配的位置中说话。它并未从周的政治语言之外另起一套天下，反而是在同一套语言中争夺谁有资格居于最高处。

这种争夺很快落到道路而非称号上。前657年，淮泗之间的徐人取舒，\eventanchor{SA-0657-XU-SHU}{春秋·徐人取舒}提醒人们：楚北面的世界并不是齐、晋、楚三强之间一块静止的空地。相邻小国也会兼并、移动和重组。次年，齐率诸侯侵蔡，\eventanchor{SA-0656-CAI-CAMPAIGN}{春秋·齐军侵蔡}；蔡地由此成为北方联盟逼近楚境的通道，而不只是地图上的一个附属国。齐军在\eventref{SA-0656-SHAOLING}{春秋·召陵之盟}前后以周礼贡赋的语言诘问楚，这场交涉没有把楚变成齐的属国，却显示楚已不能只凭江汉纵深回避中原秩序。北方联盟能否经过蔡地、楚能否令蔡继续站在自己一边，都是比盟辞更持久的考验。

前632年的\eventref{SA-0632-CHENGPU}{春秋·城濮之战}使晋取得一段可组织诸侯的优势。就楚而言，败绩并不等于南方国家忽然失去土地、城邑或动员能力；它更具体地暴露出，若要同晋在中原竞争，楚必须使北上的路线不只通到一次战场，还要通到能够供给军队、容纳使者并约束属国的节点。前634年楚灭夔，\eventanchor{SA-0634-CHU-KUI}{春秋·楚灭夔}。夔的确切地望和楚军行动范围仍待历史地理考订，但经传所保留的灭国年次足以说明，楚同时在汉水上游和中原方向经营关系。把夔仅仅视作一块待取的土地，会遮住它所处通道、属从与出兵资格的意义。

楚的王位同样不是一条无需维护的继承线。前626年，世子商臣以宫廷武力取代成王，成为穆王，\eventanchor{SA-0626-CHU-MU-ASCEND}{春秋·楚穆王即位}。经年可以固定这次非常继承，围宫与临终请求一类细节则主要属于《左传》展开的叙事，不能补作可核的宫廷纪录。较可把握的是，太子、宫卫和执政集团的站队足以改变决策中心；楚向北的政策并非由一个不变的“国家意志”自行延续。

穆王时期的连续兼并，让江汉腹地的政治边界变得更实。前623年，楚围江，江降而亡，\eventanchor{SA-0623-CHU-JIANG}{春秋·楚灭江}；次年又围六、灭六，\eventanchor{SA-0622-CHU-LIU}{春秋·楚灭六}。江、六故地、人口去向和实际控制范围都不能由几行经传精确画出，但这两件事表明齐霸衰落以后，原先可以期待北方援助的淮汉小国，已更难维持原有位置。楚取得的不是一张现代意义的疆界图，而是若干必须重新安置、守卫和调用的城邑关系。征服扩大了可用资源，也扩大了中心需要承担的供给、安抚与防叛成本。

前614年，楚庄王即位，\eventanchor{SA-0614-CHU-ZHUANG-ASCEND}{春秋·楚庄王即位}。后世常以“不鸣则已”的整齐故事塑造他的早年；这类人物箴言不能替代一份楚廷权力分配表。庄王首先面对的是继承之后如何让令尹、公族和各地军事资源重新服从共同指挥。前611年楚伐庸并灭其国，\eventanchor{SA-0611-CHU-YONG}{春秋·楚灭庸}。庸的范围以及联合反楚者的组成，主要须依赖《左传》叙事和区域研究讨论；可以确定的是，上游通路一旦脱离控制，楚的北进便没有稳固后方。此役的意义不在于替庄王预写“霸业”，而在于恢复了江汉上游可以出兵、通行和征调的条件。

这个条件仍须同国内强族讨价还价。前605年，若敖氏斗越椒起兵，庄王平定其族，\eventanchor{SA-0605-CHU-RUOAO}{春秋·若敖氏之乱}。预言、品评和人物性格是《左传》最富张力的部分，却不能证明一场冲突只由一种私人动机造成。更稳妥的理解是：控制军队、封邑和婚姻网络的贵族，并非国君命令的透明传声筒；庄王在胜利后重新分配令尹与公族资源，才较有可能把对外行动维持为连续行动。次年楚军北至洛水，传统叙事所谓“问鼎”的言辞，正是在这样的力量背景上被后世叙事放大。鼎是周王权的象征，楚军北临洛水则是可系年的压力；至于庄王与王孙满逐字如何问答，应当留在《左传》《国语》塑造的政治记忆中，而不能当作现场速记。

前598年，楚再围郑，晋的救援未能阻止郑暂时服楚，\eventanchor{SA-0598-CHU-ZHENG-SIEGE}{春秋·楚围郑}。围城显示的并不只是楚军“强”或晋军“弱”。郑处在晋楚通路之间，城中粮食、守军和卿大夫究竟能支撑多久，关系到它能否承受两大国各自提出的从属条件；晋则须把盟友、行军和军粮在更远距离上协调起来。楚能让郑转向，削弱的是晋在中原的关系网，而不只是取得一次城下胜负。

\begin{event}{公元前597年}{邲之战与楚的北方声势}{可靠纪年}{邲}
\eventref{SA-0597-BI}{春秋·邲之战}把这种压力推到会战。楚军击败晋军，直接结果是郑、陈、蔡等国不能再把晋的保护视为一项无需检验的保证；中期影响则是晋楚双方都要重新以出兵、赏赐和会盟争取这些国家。胜利并没有使楚取得可以向中原直接征税的“天下”。楚军离开江汉越远，越依赖沿路城邑、属国与将领的合作；晋的失败也没有消灭其北方腹地和动员能力。

《春秋》给出交战与败绩的简短骨架。《左传》把晋军进退、将领分歧和楚庄王的处置写成一场极具伦理张力的战争，《国语》保存不同的国别政治语言，《史记·楚世家》则把它纳入楚国盛衰的汉代长线。它们共同使人看见战争与联盟如何相互牵动，却不足以复原军队的精确人数、每一次军令或任何一方唯一的意图。邲之战打破的不是一个永远不动的晋霸，而是“北方秩序已经定局”的假定。
\end{event}

约前591至前590年，庄王卒、共王继位，\eventanchor{SA-0590-CHU-ZHUANG-DEATH}{春秋·楚庄王卒与共王继位}。新君接手的是一张尚可北进、却须不断付费维持的关系网。前571年，右司马公子申被杀，\eventanchor{SA-0571-CHU-GONGZI-SHEN}{春秋·楚杀公子申}；前570年共王卒、康王继位，\eventanchor{SA-0570-CHU-GONG-ASCEND}{春秋·楚康王继位}；前568年又杀公子壬夫，\eventanchor{SA-0568-CHU-RENFU-KILLED}{春秋·楚杀公子壬夫}。传文中关于贿赂、罪状和个人品行的解释，往往带有处置完成后的胜方视角；但连续的处死与继承至少说明，对外盟网带来的职位、馈赠和军事资源，同样会成为楚廷内部争夺的对象。

郑的选择正被这种内部与外部的牵连反复改变。前565年子囊伐郑，\eventanchor{SA-0565-CHU-ATTACK-ZHENG}{春秋·楚子囊伐郑}；次年楚又再伐郑，\eventanchor{SA-0564-CHU-SECOND-ATTACK-ZHENG}{春秋·楚再伐郑}。郑不是一枚由晋楚任意推动的棋子：它会计算城邑、民众和眼前军事压力，楚也必须把保护属国的承诺同己方行军、补给相配合。故而郑一时向楚倾斜，不应被写成一种永久归属；楚两次用兵，也不能自动证明楚廷拥有一套无摩擦的统一指挥。后来\eventref{SA-0546-MIBING}{春秋·弭兵之会}承认晋楚各自的从属网络，所冻结的是两强直接决战的一部分压力，并未消除这些城邑在两边之间反复转换的成本。

康王在前545年卒，郏敖继位，\eventanchor{SA-0545-CHU-KANG-DEATH}{春秋·楚康王卒与郏敖继位}。短促的继承没有消解宗室竞争。前541年公子围掌权，后称灵王，\eventanchor{SA-0541-CHU-LING-ASCEND}{春秋·楚灵王即位}；同年公子比出奔晋，\eventanchor{SA-0541-CHU-ZIBI-EXILE}{春秋·楚公子比出奔晋}。出奔者在异国能获得何种庇护，材料不足以列成明确的契约；但这件事使楚的王位问题越出郢都，也为日后政变保留了一条外部关系线。灵王能调动军政资源，并不等于他已使宗室和封君接受了无可争议的继承秩序。

前538年，楚在申会集诸侯及淮夷，\eventanchor{SA-0538-SHEN-CONFERENCE}{春秋·楚会诸侯于申}，又伐吴、灭赖，\eventanchor{SA-0538-CHU-LAI}{春秋·楚灭赖}。参会者的名单不是一张各方同等服从的贡赋簿，盟辞也没有完整保存；它所显示的，是楚灵王试图把陈、蔡、郑、许、徐和淮上群体置入可调用的网络。会盟的好处是能在一处集中展示军力与名分，代价则是属国必须承担出兵、供给和服从的压力。灵王随后的前537年杀屈申，\eventanchor{SA-0537-CHU-QU-SHEN}{春秋·楚杀屈申}，不宜被简单解释为淮上扩张的直接后果，却提醒人们：对外动员与内部清洗可以同时发生，而非彼此自然证明。

前534年，楚灭陈，\eventanchor{SA-0534-CHU-CHEN}{春秋·楚灭陈}，又杀陈国行人干征师，\eventanchor{SA-0534-CHU-KILLS-CHEN-ENVOY}{春秋·楚杀陈使}；前531年，蔡灵公被杀，蔡亦一度被灭，\eventanchor{SA-0531-CHU-CAI}{春秋·楚灭蔡}。陈、蔡内部的继承纠纷和对楚的回应，是《左传》解释楚军介入的主要线索，不能浓缩为楚君的一项单独欲望。较清楚的政治结果是，楚把原先须经盟约和当地精英维系的影响，改为更直接的强制控制。这样的控制能迅速清除不确定性，也会使原有地方网络失去继续替楚传递命令的理由。

\begin{debate}[灭国与复国之间的统治成本]
楚灵王前529年失国而死，\eventanchor{SA-0529-CHU-LING}{春秋·楚灵王失国}；同年陈侯、蔡侯得以归国，\eventanchor{SA-0529-CHEN-CAI-RESTORATION}{春秋·陈蔡复国}。弃疾继位后成为平王，次年前后开始重组王室与令尹集团，\eventanchor{SA-0528-CHU-PING-ASCEND}{春秋·楚平王重整楚廷}。传世材料对于政变的参与者、灵王末路和复陈蔡的用意各有鲜明叙事，不能把它们归结为平王忽然“仁厚”，也不能武断认定为单一的精算策略。陈蔡复国所能较稳地说明的是：撤去一国君主并不等于立即获得一套低成本的地方行政；楚廷若要恢复命令，还须重新同当地的宗族、城邑和边地关系协商。
\end{debate}

前527年楚诱杀戎蛮子，\eventanchor{SA-0527-CHU-RONGMAN}{春秋·楚杀戎蛮子}。经文所存的行动可以系年，诱杀场景与言辞则带有叙事的伦理色彩；“戎蛮”也不应被当作固定民族实体。此事更适宜放回边地控制的难题中理解：楚试图压住首领与通行区域，所付出的可能是信誉和长期关系的代价，而非得到一条从此绝对安全的边界。

吴军在\eventref{SA-0506-BOJU}{春秋·柏举之战}攻入郢都，楚昭王出奔，给这种治理结构作了一次严酷检验。郢都失守当然是中心遭到击穿，却不等于楚的江汉腹地、地方关系和可动员人力同时消失。前505年秦军援楚，帮助楚逐步收复核心地区；申包胥哭秦庭的完整场景见于《左传》《国语》与《史记》的文学化叙事，不能据以复原求援过程。较能确定的是，秦的援助、楚国仍在运作的腹地，以及吴军难以长期覆盖陌生而辽阔的区域，共同限制了柏举胜利转化为持久占领。同年越军乘吴主力仍在楚地而入吴，也使远征者不得不重新计算后方风险。

楚恢复后没有回到一切如旧的状态。前496年，楚公子结灭顿，\eventanchor{SA-0496-CHU-DESTROYS-DUN}{春秋·楚灭顿}；前495年又灭胡，\eventanchor{SA-0495-CHU-HU}{春秋·楚灭胡}。这些小国的准确位置和灭国后人口去向难以从传世材料重建，然而连续清除淮汉间的缓冲政治体，显示楚仍在修补吴楚竞争所暴露的东部通道。前494年楚联陈、随、许围蔡，\eventanchor{SA-0494-CHU-CAI-SIEGE}{春秋·楚围蔡}；它同吴越战事是否构成一项统一战略，现有材料不能裁定，但蔡在楚、吴之间转换依附，已足以使楚把围城当作恢复影响的一种办法。

东部并未因楚的恢复而安定。前486年楚伐陈，吴出兵救陈，\eventanchor{SA-0486-CHU-CHEN-CAMPAIGN}{春秋·楚伐陈，吴救陈}；前482年楚公子申再伐陈，\eventanchor{SA-0482-CHU-ATTACK-CHEN}{春秋·楚再伐陈}。经文能够固定战争与救援的年次，却不保存吴军规模、到达时点或楚军全部目的。陈所在的通道位置，使它既能成为楚维持淮上影响的对象，也能成为吴北上展示保护能力的支点。楚在柏举后仍可用兵，不能证明它已经解决了所有边防问题；恰恰相反，这些反复的行动说明，恢复国家并不等于恢复一条无需反复征发、交涉和守卫的稳定边界。

\begin{sourcenote}[江汉国家怎样进入传世材料]
楚的早期形成尤其需要区分材料层次。约前704年的称王不在《春秋》经年的覆盖内，主要由《史记·楚世家》及相关传统构成国别长线；前634年以后，灭国、会盟、围城、弑君和复国多可由《春秋》系年，但经文通常只留下行动者与结果。《左传》最能展示庄王、灵王、陈蔡与吴楚战争的事件次序和政治语言，却不是军令簿、税册或谈判录；《国语·楚语》保存国别说辞传统，也不能逐句还原朝廷发言。《史记》在汉代将分散材料编为楚国盛衰，更应明示为后设整理。区域考古与历史地理研究可帮助追问江汉、上游山道与淮域城邑如何约束动员，却不能替失传材料回答每一处城邑的赋役、人口或每位贵族的动机。所谓“楚是蛮夷”首先是部分传世书写的分类法，不能取代对楚自身资源、精英网络与政治选择的解释。
\end{sourcenote}

\begin{turningpoint}[制度得失：纵深必须被组织起来]
楚以江汉腹地、河流通道、兼并城邑和北方属国建立了持续竞争的条件。称王、问鼎、围郑与邲之胜解决的，是如何让南方大国在周王室名分和中原会盟中取得发言权；代价则是每一次北进或东控，都要由令尹、公族、封君和地方精英把命令转成粮秣、兵员与道路。若敖氏之乱、灵王的扩张与失国、陈蔡复国，以及柏举后的重建，反复显出同一限度：空间可以提供恢复的余地，却不会自行提供稳定的征发和服从。楚没有因郢都失守而消失，也没有因收复郢都便消除这种负担。
\end{turningpoint}
